\chapter{Existence of the families and the Galois-representation inputs}\label{ch:fam}

This chapter covers two parts of the archive. The first reduces the statement \texttt{TopG}, which says that for every host requirement there are connected families of packages at all small $\eps$, to a hypothesis record \texttt{HypsAlt} and the parametrisation \texttt{Green} of cuspidal characters of $\GL_d(\F_g)$; that record is then produced from named theta-correspondence inputs, a covolume bound and the statement \texttt{RecTExists}. The second part states exactly what the two ``literature'' inputs threaded to the top say: \texttt{RecTExists}, a weak local Langlands map, and \texttt{Row0b}, a truncated local--global compatibility for regular algebraic polarized cuspidal representations over CM fields. Both are proved inside the archive and are plugged into the final theorem through the record \texttt{OpenLiteratureStatements}. The Sha bound and the Chebyshev inequality used next to these results are treated elsewhere.

\section{From a hypothesis record to connected families}

Throughout, a \emph{class of families} is a predicate on families $\mathrm{Fm}$ of packages at a parameter $\eps$ (the archive's \texttt{FamP eps}); a family has a set $\mathrm{Win}$ of scales $x$ and, for each large $x$, finitely many packages $D(x,i)$ with size data $L(x,i)$. A property $P(\mathrm{Fm},C_H)$ is \emph{closed} on a class $\mathcal C$ if there is $\eps_1>0$ such that for every $0<\eps\le\eps_1$ there is a threshold $K$ with $P(\mathrm{Fm},C_H)$ for every $\mathrm{Fm}\in\mathcal C$ at $\eps$ and every $C_H$ for which $\mathrm{Fm}$ is large for $(K,C_H)$.





\begin{theorem}[Connected families from closed links and a design]\label{fam:top-des-win}
\lean{UpperTheory.Top.topDesWin}
% uses-lean: Formalism.MainNotions.Closed, Formalism.MainNotions.Des, Formalism.MainNotions.DesignExists, Formalism.MainNotions.LinkN, Formalism.MainNotions.LinkEta, Formalism.MainNotions.LinkIRR, Formalism.MainNotions.LinkDIVt, Formalism.MainNotions.LinkG, Formalism.MainNotions.LinkS, Formalism.MainNotions.LinkGLOBAL, UpperTheory.Top.HypV
\uses{geob:closed, geob:lowerlink, geob:upperlink-fam} % resolved
\uses{geob:topg}
Let $\mathcal C,\mathcal C',\mathcal E$ be classes of families and $C_{\mathrm{abs}},C_0,e_D,e_L$ real numbers. Assume:
\begin{enumerate}
\item a design exists in the class of families of $\mathcal C'$ with good twists and smooth $D_{\Gamma}$, with constant $C_0$ (\texttt{DesignExists}): there is $\eps_0>0$ such that for every $0<\eps\le\eps_0$ and every threshold $K$ there are a family $\mathrm{Fm}$ in that class and $C_H\ge0$ with $\mathrm{Fm}$ large for $(K,C_H)$, chosen with constants $(C_0,C_H)$, and satisfying \texttt{TuSize} and \texttt{DGamSize};
\item $\mathcal C'\subseteq\mathcal C$ and $\mathcal C'\subseteq\mathcal E$, and every family of $\mathcal E$ satisfies \texttt{HypV};
\item on the class $\mathrm{Des}(\mathcal C,C_0)$ the following are closed: \texttt{LinkN}; \texttt{LinkEta}; \texttt{LinkIRR}; the rank equality \texttt{SmallLevelRankEq}; ``\texttt{TuSize} and \texttt{DGamSize} imply \texttt{LinkDIVt} with $(C_{\mathrm{abs}},C_H)$''; ``being chosen with $(C_0,C_H)$ implies full rank of $\nu$ and \texttt{MuRank}''; \texttt{LinkG}; ``\texttt{TuSize} implies \texttt{LinkS} with $C_0$''; and ``\texttt{TuSize} and \texttt{DGamSize} imply \texttt{LinkGLOBAL} with $(C_{\mathrm{abs}},C_H,C_0,e_D,e_L)$''.
\end{enumerate}
Then $\mathcal E$ is connected.
\end{theorem}
\begin{proof}
% uses-lean: UpperTheory.Top.divinN_of_hypV, Formalism.MainNotions.J.linkMu_of_linkDIV, Formalism.MainNotions.muTu_of_tuSize, Formalism.MainNotions.J.muDGam_of_fact
Take the conjunction of the nine closed properties; it is closed with some $\eps_1$. Let $\eps_3$ be the constant of the design and put $\eps_1'=\min(\eps_1,\eps_3)$; the constants $e_D,e_L$ are those given. For $0<\eps<\eps_1'$, the closed properties give a threshold $K$; the design, applied to this $K$, gives a family $\mathrm{Fm}$ and $C_H$ with $\mathrm{Fm}$ large for $(K,C_H)$. So the threshold is chosen before the family. The family lies in $\mathrm{Des}(\mathcal C,C_0)$, so all nine properties hold for it. From \texttt{HypV} (which holds since $\mathrm{Fm}\in\mathcal E$), \texttt{LinkN} and \texttt{SmallLevelRankEq} the archive's lemma \texttt{divinN\_of\_hypV} gives \texttt{DIVinN}; with \texttt{LinkEta} this feeds \texttt{LinkDIVt} and gives \texttt{LinkDIV}. Then \texttt{LinkDIV}, the rank statements, \texttt{TuSize} and \texttt{DGamSize} give \texttt{LinkMu}. Now \texttt{LinkG}, \texttt{LinkS}, \texttt{LinkIRR}, \texttt{LinkDIV} and \texttt{LinkMu} are the global theorem inputs, and \texttt{LinkGLOBAL} turns them into the tied chain at every large $x\in\mathrm{Win}$.
\end{proof}

\begin{remark}[Order of quantifiers]\label{fam:rem-quantifiers}
In Theorem~\ref{fam:top-des-win} the threshold $K$ comes from the closed properties and the family is produced afterwards, large for that $K$. So the conclusion is a genuine existence statement for each small $\eps$. An empty class cannot be connected.
\end{remark}

\begin{definition}[Green's parametrisation]\label{fam:green}
\lean{GlobalTheory.GLq.Green, GlobalTheory.GLq.LabelDatum, GlobalTheory.GLq.IsPrimeLab, GlobalTheory.GLq.IsCuspOf}
For a prime $g$, a \emph{label datum} $L$ consists of an injective-on-nonzero-elements Teichm\"uller map $\overline{\F}_g\to\CC$ (multiplicative) and elements $\gamma_d\in\overline{\F}_g$ of order $g^d-1$, compatible under norms. A polynomial $P\in\F_g[X]$ is a \emph{prime label of degree $d$} if it is monic, irreducible, of degree $d$, with $P(0)\ne0$. \texttt{Green} is the statement: for every prime $g$, every label datum $L$ and every $d>0$ there is a map $\tau$ from $\F_g[X]$ to functions $\GL_d(\F_g)\to\CC$ such that
\begin{enumerate}
\item for every prime label $P$ of degree $d$, $\tau(P)$ is cuspidal with label $P$ for $L$ (its value at an element whose characteristic polynomial is irreducible with root $z$ is $(-1)^{d-1}\sum_{i<d}\theta_x(z^{g^i})$, where $\theta_x$ is the character attached by $L$ to a root $x$ of $P$), and $\tau(P)(1)=\prod_{i=1}^{d-1}(g^i-1)$;
\item $\tau$ is injective on prime labels of degree $d$;
\item every cuspidal function on $\GL_d(\F_g)$ equals $\tau(P)$ for some prime label $P$ of degree $d$.
\end{enumerate}
The predicate ``cuspidal'' (\texttt{IsCusp}) was not transcribed.
\end{definition}

\begin{definition}[Design supply]\label{fam:exs}
\lean{UpperTheory.Top.ExS}
% uses-lean: Formalism.MainNotions.DesignChoices, GlobalTheory.EX.HostExW, UpperTheory.Top.FReq.cls
\uses{host:host-exw} % resolved
For a class $\mathcal A$ and a function $c_T$, \texttt{ExS}$(\mathcal A,c_T)$ is the statement: for every $\lambda>1$ and every host requirement $G$ with \texttt{HostExW}$(\lambda,G)$, design choices with constant $3$ and $c_T$ exist in the class of families of $\mathcal A\cap G.\mathrm{cls}(\lambda)$ with good twists (\texttt{DesignChoices}).
\end{definition}

\begin{definition}[The shape of the top argument]\label{fam:top-shape-alt}
\lean{GlobalTheory.Fc.TopShapeAlt}
% uses-lean: Formalism.MainNotions.Closed, Formalism.L2.Top3.D3, Formalism.MainNotions.MinkDiv, Formalism.MainNotions.Duty
\uses{geob:closed} % resolved
\uses{fam:exs}
For propositions $V,N_R,G$ and a predicate $N$ on classes, \texttt{TopShapeAlt}$(V,N_R,N,G)$ is the statement that $G$ holds whenever one is given, for some class $\mathcal A$:
\begin{enumerate}
\item a function $c_T$ with ``Minkowski's divisibility bound implies \texttt{ExS}$(\mathcal A,c_T)$'';
\item a twist selection for $\mathcal A$ and, built on it, spectral, base-level, definite-height, column and bundle selections, and the duties \texttt{Duty} of every member of $\mathcal A$;
\item about forty statements on the selected data: bounded and primitive indefinite realisations, twist ties, row identities and divisibility of rows, the size statement \texttt{TuSize}, the class number formula, the thin-set statement \texttt{LitThin}, Wu decay, finite-dimensionality of rows, a frame statement, counting constants, theorem S on selected data, a point count, mass and volume statements, image and surjectivity statements for level projections, integrality of twisted points, periods and columns, a family-size bound, the real-place second-layer conditions, and the mass ratio and genus mass at $2$; most of these are required to be closed on the class $D_3(3,\mathcal A)$ (the design class of $\mathcal A$ with constant $3$);
\item a proof of $V$, a proof of $N_R$, and $N(D_3(3,\mathcal A))$;
\item a constant $c_I>0$ and a closed label-count statement used for \texttt{LinkIRR}.
\end{enumerate}
The arguments of item 3 are not transcribed.
\end{definition}

\begin{definition}[The hypothesis record]\label{fam:hyps-alt}
\lean{GlobalTheory.Fc.HypsAlt}
\uses{fam:top-shape-alt, fam:exs}
For a predicate $N$ on classes and a definite-side condition $D_R$, a record \texttt{HypsAlt}$(N,D_R)$ consists of a class $\mathcal A$, a raw selection $\sigma$ for $\mathcal A$, a function $c_T$, and fields matching the arguments of \texttt{TopShapeAlt} other than the selections (derived from $\sigma$), $V$ and $N_R$; the row field \texttt{rRows} is stated with $D_R$, and the last field is $N(D_3(3,\mathcal A))$. It is a structure of type \texttt{Type 1}, so its use as a hypothesis is through \texttt{Nonempty}.
\end{definition}

\begin{lemma}[A record instantiates the shape]\label{fam:goal-of-shape}
\lean{GlobalTheory.Fc.goal_of_shape_of_hyps, GlobalTheory.Fc.forall_goal_of_shape_of_hyps}
% uses-lean: GlobalTheory.Fc.DefiniteSideCondition
\uses{fam:top-shape-alt, fam:hyps-alt}
Let $V,N_R,G$ be propositions and $N$ a predicate on classes. Assume \texttt{TopShapeAlt}$(V,N_R,N,G)$, $V$, $N_R$, and a record \texttt{HypsAlt}$(N,\texttt{DefiniteSideCondition})$. Then $G$.
\end{lemma}
\begin{proof}
% uses-lean: GlobalTheory.Fc.familySelection, GlobalTheory.Fc.selectionBundleOfRawSelection, GlobalTheory.Cp.divRows3G_iff, GlobalTheory.Cp.serv_t_iff_formalism_serv_t
Apply the shape to the fields of the record. The selections are read off from the raw selection; two fields are converted by equivalences of formulations (divisibility of rows and the served-selection condition), and the row field is repackaged by forgetting one component.
\end{proof}

\begin{lemma}[A design from the record]\label{fam:design-exists}
\lean{GlobalTheory.Top.designExists_of_tuSize_aux4n_litThin}
% uses-lean: Formalism.MainNotions.DesignExists, Formalism.MainNotions.DesignChoices, Formalism.MainNotions.MinkDiv
Let $\mathcal A$ be a class and $c_T$ a function. Assume that Minkowski's divisibility bound implies design choices with constant $3$ in the class of $\mathcal A$ with good twists; that a twist selection for $\mathcal A$ is given and every member of $\mathcal A$ satisfies its duties; that ``$c_T(\eps,C_{YL})\le C_{Tu}$ implies \texttt{TuSize}'' and the class number formula are closed on the class of $\mathcal A$ with chosen data, smooth $D_{\Gamma}$ and good twists; and \texttt{LitThin}. Then a design with constant $3$ exists in that last class.
\end{lemma}
\begin{proof}
Only the top of the proof was read: the duties give aligned linked families, hence second-layer families, from which the design is assembled; the proof was not expanded further.
\end{proof}

\begin{theorem}[The shape holds for \texttt{TopG}]\label{fam:shape-closed}
\lean{UpperTheory.Top.shape_closed_link_n_of_green}
\uses{fam:green, fam:top-shape-alt, geob:topg, fam:top-des-win, fam:design-exists}
Assume \texttt{Green}. Then \texttt{TopShapeAlt}$(\mathrm{True},\mathrm{True},N,\texttt{TopG})$ holds, where $N(\mathcal C)$ is the statement that \texttt{LinkN} is closed on $\mathcal C$.
\end{theorem}
\begin{proof}
\uses{fam:top-des-win, fam:design-exists}
% uses-lean: GlobalTheory.Top.link_gl, GlobalTheory.Top.link_brest_alt, GlobalTheory.Top.closed_family_admissible_twist_link_alt, GlobalTheory.Top.selection_closed_of_wu_decay_alt, GlobalTheory.Top.wire_sel_alt, Formalism.L2.DIV.closed_linkDIVt_of_closed_bridge_kernel, GlobalTheory.Top.closed_linkEta_of_duty, GlobalTheory.Top.Irr.closed_linkIRR_of_join_count, GlobalTheory.Top.link_alt_alt_alt, GlobalTheory.Top.closed_linkG_of_duty, GlobalTheory.Top.link_alt_alt_bis, Formalism.MainNotions.DIV.closed_Csz_ge_fn
Fix the arguments of the shape and $\lambda>1$, $G$ with \texttt{HostExW}$(\lambda,G)$ and \texttt{HypV} on $G.\mathrm{cls}(\lambda)$. From the duties, the real-place conditions and the mass statements one obtains \texttt{LinkGLOBAL} and the closed rank statement. The bridge-row statements, the admissible-twist link and the Wu-decay selection combine into a closed bridge-kernel bound, which gives the closed \texttt{LinkDIVt}. The design of Lemma~\ref{fam:design-exists} is applied to the class $\mathcal A\cap G.\mathrm{cls}(\lambda)$, using \texttt{ExS} for this $(\lambda,G)$. Theorem~\ref{fam:top-des-win} is then applied with $\mathcal C$ the class of $\mathcal A$ with chosen data, $\mathcal C'$ the class of $\mathcal A\cap G.\mathrm{cls}(\lambda)$ with chosen data, $\mathcal E=G.\mathrm{cls}(\lambda)$ and constants $(C_{\mathrm{abs}},C_0,e_D,e_L)=(2,3,2,0)$. Its closed inputs are: \texttt{LinkN} (the argument $N$), \texttt{LinkEta} and \texttt{LinkG} from the duties, \texttt{LinkIRR} from \texttt{Green} and the label count, the rank equality from surjectivity of the level projections, \texttt{LinkS} from the counting statements, and \texttt{LinkDIVt} and \texttt{LinkGLOBAL} as above. The two arguments $V$ and $N_R$ of the shape are not used.
\end{proof}



\begin{remark}[The literal truths]\label{fam:rem-trivial}
The two arguments \texttt{trivial} in Theorem~\ref{fin:every-green} fill the parameters $V$ and $N_R$ of \texttt{TopShapeAlt}, which Theorem~\ref{fam:shape-closed} instantiates to \texttt{True} and never uses. They discharge no mathematical hypothesis.
\end{remark}

\begin{proposition}[A record from theta inputs, covolume and \texttt{RecTExists}]\label{fam:nonempty-hyps}
\lean{UpperTheory.Lit.nonempty_hyps_definite_side_of_closed_bound}
% uses-lean: GlobalTheory.Fc.S0.UnitaryGroupThirdPlaceHypotheses, GlobalTheory.FunPr.ThetaLocalInputs, GlobalTheory.IndefiniteRowSideFullHecke.WeilTensorModels.ThetaIntegralsAgree, GlobalTheory.IndefiniteRowSideFullHecke.WeilTensorModels.LocalMatching, GlobalTheory.CovTop.CovolumeBoundOfX, GlobalTheory.Fc.S0.NL
\uses{geob:covolbound} % resolved
\uses{fam:hyps-alt, fam:rect-exists}
Assume the records \texttt{UnitaryGroupThirdPlaceHypotheses} and \texttt{ThetaLocalInputs}, the statements \texttt{ThetaIntegralsAgree} and \texttt{LocalMatching}, the covolume bound \texttt{CovolumeBoundOfX}, and \texttt{RecTExists}; and assume that for every $c_{\mathrm{Dep}}$ the parameters of dependence built from \texttt{RecTExists} and $c_{\mathrm{Dep}}$ are linked for the condition that the bound at the places $2$ and $3$ is closed. Then there is a record \texttt{HypsAlt}$(N,\texttt{DefiniteSideCondition})$, with $N(\mathcal C)$ the statement that \texttt{LinkN} is closed on $\mathcal C$.
\end{proposition}
\begin{proof}
% uses-lean: UpperTheory.Lit.functionalsSeparate_of_howeUnram_of_siegelWeil, GlobalTheory.Mem.LiteralAutomorphy.rowRatW_of_rowTest, GlobalTheory.Glob.CovolRow.linked_bundles_l3_of_open_alt, GlobalTheory.Fc.S0.StageU.hypotheses_nonempty_of_unramified_exponents_alt
Read at the top level only. Unramified Howe duality and the Siegel--Weil formula separate functionals; with the test-vector assembly this gives rationality of the row functions; a Fourier-coefficient fit and an evaluation grid follow from the theta inputs; row functions on the grid come from discrete series and Casselman--Wallach. The parameters of dependence are opened and linked, with the geometry predicate equal to the constant \texttt{True} (\texttt{anyGeometry}); the covolume bound gives the linked bundles. All of this is fed to \texttt{hypotheses\_nonempty\_of\_unramified\_exponents\_alt}, which produces the record. Its proof, which must construct packages at every scale, lives in \arch{Development/AutomorphicArithmetic/ThetaCorrespondence/ThetaPairLifts/ThirdPlaceSeries/} and was not expanded.
\end{proof}



\begin{definition}[The older copy of the shape and the record]\label{fam:top-shape-old}
\lean{GlobalTheory.Fc.TopShape, GlobalTheory.Fc.Hyps}
\uses{fam:top-shape-alt, fam:hyps-alt}
\texttt{TopShape} and \texttt{Hyps} are earlier versions of Definitions~\ref{fam:top-shape-alt} and~\ref{fam:hyps-alt}, with the same argument and field names in the same order; their types were not compared field by field.
\end{definition}

\begin{theorem}[The older top chain]\label{fam:shape-win}
% overlaps (deduplicated): fin:plug
\lean{UpperTheory.Top.shapeWin, GlobalTheory.Fc.shape_of, UpperTheory.Lit.hypsE₂}
% uses-lean: GlobalTheory.Published.LocalLanglandsGL, UpperTheory.Top.zeroFreeAll_of_topG
\uses{an:zerofreeall-topg, fam:llc} % resolved
\uses{fam:top-shape-old, fam:green, geob:topg, fam:top-des-win, fam:llc}
Assume \texttt{Green}. Then \texttt{TopShape}$(\mathrm{True},\mathrm{True},N,\texttt{TopG})$ holds (\texttt{shapeWin}); a \texttt{Hyps} record instantiates \texttt{TopShape} (\texttt{shape\_of}); with \texttt{VWin} and a \texttt{Hyps} record one gets \texttt{ZeroFreeAll} and every character (\texttt{plugWin}, \texttt{plugAll}). The record is produced by \texttt{hypsE$_2$} from records \texttt{HypothesesUE}, \texttt{Hypotheses$_2$}, the theta statements, a covolume statement \texttt{CovolX}, \texttt{LocalLanglandsGL}, and the statement that for every $c_{\mathrm{Dep}}\in\R$ every closed family at the bundle points of the parameters \texttt{prm1}$(c_{\mathrm{Dep}})$ satisfies the bound at the places $2$ and $3$ (\texttt{BoundAtTwoAndThreeForFamily}; inside \texttt{everyE$_2$} this is supplied by \texttt{boundAtTwoThree\_at\_prm\_of\_prints}), and \texttt{everyE$_2$} combines these.
\end{theorem}
\begin{proof}
\uses{fam:top-des-win}
The proof of \texttt{shapeWin} has the same structure as that of Theorem~\ref{fam:shape-closed}, with the older link lemmas (\texttt{link\_ex}, \texttt{link\_eta}, \texttt{link\_irr}, \texttt{link\_g}, \texttt{link\_alt}) and \texttt{topDesWin} at the end. The other statements are one-line compositions.
\end{proof}

\begin{remark}[The older chain is not used]\label{fam:rem-dead-copy}
The final theorem reaches the existence of families through Theorems~\ref{fin:theta-local}, \ref{fin:every-green} and~\ref{fam:shape-closed}. The older chain of Theorem~\ref{fam:shape-win} is not referenced by it: the only use of \texttt{plugAll} is in \texttt{everyE$_2$}, and no file other than its own mentions \texttt{everyE$_2$}. In particular the hypothesis \texttt{LocalLanglandsGL} of \texttt{hypsE$_2$} is not an input of the final theorem.
\end{remark}

\section{The weak local Langlands map}

Fix a number field $E$, an algebraic closure $\Omega$, a valuation subring $A$ of $\Omega$ lying over a finite place $w$ of $E$; this fixes a decomposition group and a Weil group $W_w$. Write $\mathrm{Irr}$ for the set of pairs $(n,\pi)$ with $\pi$ an irreducible smooth representation of $\GL_n(E_w)$ up to isomorphism, and $\mathrm{WD}$ for isomorphism classes of Weil--Deligne representations of $W_w$.

\begin{definition}[Good representations]\label{fam:good}
\lean{UpperTheory.Zero.LLC.GoodCusp, UpperTheory.Zero.LLC.Good}
An element $y\in\mathrm{Irr}$ is a \emph{good cuspidal} if either $y$ is a representation of $\GL_1$, or $E_w$ is isomorphic to $\Q_g$ for a prime $g$ (by a continuous ring isomorphism) and $y$ is a twist by a smooth character of some $y_0$ having nonzero vectors fixed by the first congruence subgroup $K_1$ (that is, $y_0$ has depth zero). An element $x\in\mathrm{Irr}$ is \emph{good} if it has a cuspidal support all of whose members are good cuspidals.
\end{definition}

\begin{definition}[The type law]\label{fam:type-law}
\lean{UpperTheory.Zero.LLC.TypeLaw}
A map $\mathrm{rc}\colon\mathrm{Irr}\to\mathrm{WD}$ satisfies the \emph{type law} if, whenever $E_w\cong\Q_g$ and $E'/E$ is an extension with a place $w'$ over $g$ that is unramified over $g$ with residue field of size $g^n$, $n>0$, and $\sigma$ is an irreducible representation of $\GL_n(E_w)$ inflating a cuspidal representation $\tau$ of $\GL_n(\F_g)$ with Green character attached to a regular character $\theta$, then $\mathrm{rc}(\sigma)$ has monodromy $0$ and its characteristic polynomial on inertia is $\prod_{j<n}(X-\theta(z^{g^j}))$ at the corresponding elements.
\end{definition}

\begin{definition}[The weak laws]\label{fam:laws-t}
\lean{UpperTheory.Zero.LLC.LawsT, UpperTheory.Zero.LLC.IsRecT}
\uses{fam:good, fam:type-law}
Let $a\colon W_w\to E_w^{\times}$ be a local Artin map. A map $\mathrm{rc}\colon\mathrm{Irr}\to\mathrm{WD}$ satisfies the \emph{weak laws} for $a$ if:
\begin{enumerate}
\item $\dim\mathrm{rc}(n,\pi)=n$, and on characters $\mathrm{rc}(\chi)=\chi\circ a$;
\item if $\pi$ is the Langlands quotient of a list of segments, $\mathrm{rc}(\pi)$ is the corresponding sum over segments of the twisted images of the first cuspidals tensored with special representations;
\item $\mathrm{rc}$ of a cuspidal is irreducible (one direction only);
\item $\mathrm{rc}(\pi\otimes|\det|^s)=\mathrm{rc}(\pi)\otimes|\cdot|^s$;
\item if $\pi$ has central character $\omega$, the absolute value of $\det\mathrm{rc}(\pi)(\sigma)$ is $|\omega(a(\sigma))|$;
\item for good $\pi$ only: $\mathrm{rc}$ respects twists by characters; $\det\mathrm{rc}(\pi)=\omega\circ a$; $\mathrm{rc}$ of the contragredient is the dual; $\mathrm{rc}(\pi')=\mathrm{rc}(\pi)$ implies $\pi'=\pi$; and depth zero implies tame;
\item the type law of Definition~\ref{fam:type-law}.
\end{enumerate}
\texttt{IsRecT}$(\mathrm{rc})$ means that $\mathrm{rc}$ satisfies the weak laws for some local Artin map $a$.
\end{definition}

\begin{definition}[Existence and uniqueness of weak maps]\label{fam:rect-exists}
\lean{UpperTheory.Zero.LLC.RecTExists, UpperTheory.Zero.LLC.RecTUnique}
\uses{fam:laws-t, fam:good}
\texttt{RecTExists} is the statement: for every $E,\Omega,A,w$ as above there is a map $\mathrm{rc}$ with \texttt{IsRecT}$(\mathrm{rc})$. \texttt{RecTUnique} is the statement: any two maps with \texttt{IsRecT} agree on good representations.
\end{definition}

\begin{definition}[Junk parameters]\label{fam:junk-ex}
\lean{UpperTheory.Zero.LLC.JunkEx, UpperTheory.Zero.LLC.CuspEx, UpperTheory.Zero.LLC.CuspT}
\texttt{JunkEx} is the statement: for all $E,\Omega,A,w$ and $n\ge2$ there is an irreducible Weil--Deligne representation $J$ of dimension $n$ with unitary determinant such that, if $E_w\cong\Q_g$, $J$ is not a twist of a tame representation. \texttt{CuspEx} is the statement that for every local Artin map $a$ there is a map $r\colon\mathrm{Irr}\to\mathrm{WD}$ with the cuspidal laws \texttt{CuspT}: irreducible and dimension-preserving on cuspidals, class field theory on characters, compatible with $|\det|^s$, the absolute value of the central character on cuspidals, twists, central character and duals on good cuspidals, injectivity on good cuspidals, a good cuspidal of depth zero going to a tame representation, and the type condition \texttt{TypeAt}$(x, r(x))$ for every cuspidal $x$.
\end{definition}

\begin{lemma}[Cuspidal maps exist]\label{fam:cusp-ex}
\lean{UpperTheory.Zero.LLC.cuspEx}
% uses-lean: UpperTheory.Zero.LLC.goodMap, UpperTheory.Zero.LLC.junkEx
\uses{fam:junk-ex, fam:good}
\texttt{CuspEx} holds.
\end{lemma}
\begin{proof}
Take a map $r_0$ on good cuspidals with the required properties (\texttt{goodMap}), a choice $J(n)$ of junk parameters (\texttt{junkEx}) and an exponent $e$. Define $r(x)=r_0(x)$ if $x$ is a good cuspidal and $r(x)=J(\dim x)\otimes|\cdot|^{e(x)}$ otherwise. A cuspidal that is not good has dimension at least $2$, so the junk value is irreducible of the right dimension; the remaining laws only concern good cuspidals or follow from the choice of $e$.
\end{proof}

\begin{theorem}[\texttt{RecTExists} holds]\label{fam:rect-exists-thm}
\lean{UpperTheory.Zero.LLC.recTExists}
% uses-lean: UpperTheory.Zero.LLC.artEx, GlobalTheory.LocLanglands.exists_isLExt, UpperTheory.Zero.LLC.Pr.exh, UpperTheory.Zero.LLC.Pr.multi, UpperTheory.Zero.LLC.extDim, UpperTheory.Zero.LLC.extTw, UpperTheory.Zero.LLC.extDual, UpperTheory.Zero.LLC.extInj, UpperTheory.Zero.LLC.extDepth, UpperTheory.Zero.LLC.typeIff
\uses{fam:rect-exists, fam:cusp-ex, fam:laws-t}
For every number field $E$, $\Omega$, $A$ and finite place $w$ there is a map $\mathrm{Irr}\to\mathrm{WD}$ satisfying the weak laws.
\end{theorem}
\begin{proof}
\uses{fam:cusp-ex}
A local Artin map $a$ exists (\texttt{artEx}). Lemma~\ref{fam:cusp-ex} gives $r$ on cuspidals. Langlands exhaustion and the multisegment classification extend $r$ to a map $\mathrm{rc}$ on all of $\mathrm{Irr}$ through segments, agreeing with $r$ on cuspidals. The laws are then checked one by one: dimension and central characters from the extension lemma, twists by $|\det|^s$ and by characters, duals, injectivity and depth on good classes from the corresponding lemmas, and the type law by reducing to cuspidals inflated from $\GL_n(\F_g)$.
\end{proof}

\begin{remark}[\texttt{RecTExists} is not local Langlands]\label{fam:rem-not-llc}
The field of \texttt{OpenLiteratureStatements} called \texttt{localLanglandsExists} has type \texttt{RecTExists}. Compared with Definition~\ref{fam:laws} it asserts no bijectivity and no $\epsilon$-factors (of pairs), no product, determinant, support, spherical or Satake laws, and only one direction of ``cuspidal iff irreducible''. Twists, central characters, duals, injectivity and depth are required only on good representations. The map built in Theorem~\ref{fam:rect-exists-thm} takes, on every cuspidal that is not good (for instance every cuspidal of positive depth, and every non-character cuspidal when $E_w$ is not $\Q_g$), the junk value $J(n)\otimes|\cdot|^{e}$, which depends only on the dimension up to the twist. Off the good classes the value is not unconstrained: the laws \texttt{dim}, \texttt{langlands} (the segment-sum formula), \texttt{cusp}, \texttt{nu}, \texttt{absCentral} and \texttt{type} still apply there. But twists, central characters, duals, injectivity and depth pin $\mathrm{rc}$ down only on representations built from good cuspidals.
\end{remark}

\begin{definition}[The full local Langlands laws]\label{fam:laws}
\lean{GlobalTheory.LocLanglands.Laws, GlobalTheory.LocLanglands.IsRec}
For a local Artin map $a$ and a family $\epsilon$ of local constants, a map $\mathrm{rc}\colon\mathrm{Irr}\to\mathrm{WD}$ satisfies the \emph{laws} if it is bijective; preserves dimension; is $\chi\mapsto\chi\circ a$ on characters; matches the $\epsilon$-factors of pairs (cuspidal of $\GL_{m+1}$ with any of $\GL_m$) with $\epsilon$ of the tensor product; respects twists, central characters and contragredients; satisfies ``$x$ cuspidal iff $\mathrm{rc}(x)$ irreducible''; is compatible with Langlands quotients of segments, with characters composed with the determinant, with products, and with cuspidal supports on characteristic polynomials; sends spherical representations exactly to unramified ones, Satake parameters to the corresponding diagonal representation, and depth zero exactly to tame. \texttt{IsRec}$(\mathrm{rc})$ means the laws hold for some $a$ and some $\epsilon$ satisfying Deligne's conditions.
\end{definition}

\begin{definition}[The local Langlands statement]\label{fam:llc}
\lean{GlobalTheory.Published.LocalLanglandsGL}
\uses{fam:laws}
\texttt{LocalLanglandsGL} is the statement: for all $E,\Omega,A,w$ there is exactly one map $\mathrm{rc}$ with \texttt{IsRec}$(\mathrm{rc})$, and any two such maps for $(E,w)$ and $(E',w')$ are compatible under every compatible isomorphism of the local data.
\end{definition}

\begin{lemma}[Full laws imply weak laws]\label{fam:isrect-of-isrec}
\lean{UpperTheory.Zero.LLC.lawsT_of_laws, UpperTheory.Zero.LLC.isRecT_of_isRec}
% uses-lean: GlobalTheory.Published.DepthZeroTypeInertia
\uses{fam:laws, fam:laws-t}
If $\mathrm{rc}$ satisfies the laws and the type law, it satisfies the weak laws. Assume the statement \texttt{DepthZeroTypeInertia}. Then \texttt{IsRec}$(\mathrm{rc})$ implies \texttt{IsRecT}$(\mathrm{rc})$.
\end{lemma}
\begin{proof}
Each weak law is a field of the laws or its restriction to good representations; injectivity comes from bijectivity. \texttt{DepthZeroTypeInertia} supplies the type law.
\end{proof}

\begin{lemma}[\texttt{LocalLanglandsGL} only conditionally]\label{fam:llgl-conditional}
\lean{GlobalTheory.T.local_langlands_of_published, GlobalTheory.LocLanglandsBridge.localLanglandsGL_of_depthZero_zelevinsky_bernstein}
% uses-lean: GlobalTheory.Published.LocalLanglandsExists
\uses{fam:llc}
Assume \texttt{LocalLanglandsExists} (existence of a map with the literature laws \texttt{IsRecLit}), the local converse theorem, depth-zero cuspidal local Langlands, depth-zero induction, a Brauer degree-zero statement, Chevalley's unit theorem, global Artin reciprocity, Langlands exhaustion, the multisegment classification, standard cosocles, Zelevinsky's polynomial ring, the geometric lemma and Bernstein--Zelevinsky finiteness. Then \texttt{LocalLanglandsGL} holds.
\end{lemma}

\begin{remark}[\texttt{LocalLanglandsGL} is never proved outright]\label{fam:rem-llgl}
A textual search for declarations and local proofs whose conclusion is \texttt{LocalLanglandsGL} finds ten; each is Lemma~\ref{fam:llgl-conditional} or applies it to fields of a hypothesis record. The same search for \texttt{LocalLanglandsExists} finds only its derivation from \texttt{LocalLanglandsGL}. So in the archive \texttt{LocalLanglandsGL} is stated and used as a hypothesis but never proved outright. On the main chain the corresponding record fields have type \texttt{RecTExists} instead.
\end{remark}

\begin{lemma}[An identity function]\label{fam:ll-of-alt}
\lean{UpperTheory.One.M.New.llOf_alt, GlobalTheory.T.recTExists_of_published_inputs}
\uses{fam:rect-exists}
Assume \texttt{RecTExists} and eleven further statements (local converse, depth-zero induction, Brauer degree zero, Chevalley units, global Artin reciprocity, Langlands exhaustion, multisegments, standard cosocles, Zelevinsky's polynomial ring, the geometric lemma, Bernstein--Zelevinsky finiteness). Then \texttt{RecTExists}.
\end{lemma}
\begin{proof}
The conclusion is the first hypothesis; the others are discarded. \texttt{recTExists\_of\_published\_inputs} is proved by \texttt{llOf\_alt}.
\end{proof}

\begin{remark}[The identity is harmless]\label{fam:rem-identity}
Lemma~\ref{fam:ll-of-alt} is on the main chain twice: through \texttt{local\_langlands\_gl\_recT\_exists\_of\_literature}, and through \texttt{recTExists\_of\_published\_inputs} inside the record produced by Proposition~\ref{fam:nonempty-hyps}. In both places its first argument is a field of type \texttt{RecTExists} that is filled by Theorem~\ref{fam:rect-exists-thm}. The lemma is not circular reasoning, and its eleven extra hypotheses carry no weight.
\end{remark}

\section{The truncated local--global compatibility}

A \emph{rec notion} is a predicate on maps $\mathrm{rc}\colon\mathrm{Irr}\to\mathrm{WD}$ for every local datum. For a Weil--Deligne representation $y$ at $w$:
$y$ is \emph{abelian on commutators} (\texttt{Tor}) if the characteristic polynomial of $y$ at every commutator $[s,t]$ of $W_w$ is $(X-1)^{\dim y}$;
$y$ is \emph{tame} (\texttt{TmW}) if its characteristic polynomial at every wild-inertia element is $(X-1)^{\dim y}$;
$y$ is \emph{square-free} (\texttt{Sqf}) if for every $\gamma$ topologically generating inertia modulo wild inertia, every root of multiplicity at least $2$ of the characteristic polynomial of $y(\gamma)$ is a $(q_w-1)$-th root of unity.

\begin{definition}[The local conditions]\label{fam:away-b}
\lean{UpperTheory.GalStar.IsQpBig, UpperTheory.GalStar.AwayB, UpperTheory.GalStar.Tor, UpperTheory.GalS.TmW, UpperTheory.GalStar.Sqf}
The place $w$ is \emph{$\Q_g$-big for $n$} if $E_w\cong\Q_g$ continuously for a prime $g$ with $n=1$ or $16n^2<g$. For Weil--Deligne representations $x,y$ at $w$, \texttt{AwayB}$(x,y)$ means: $\dim x=\dim y$; if $y$ is abelian on commutators with monodromy $N=0$, then $x=y$; and if $w$ is $\Q_g$-big for $\dim y$ and $y$ is tame, square-free, with $N=0$, then $x$ has $N=0$ and $x$, $y$ have the same characteristic polynomial at every element of degree $0$ (that is, on inertia).
\end{definition}

\begin{definition}[The statement \texttt{Row0b}]\label{fam:props-b}
\lean{UpperTheory.GalStar.AboveB, UpperTheory.GalStar.PropsB, UpperTheory.GalStar.Row0b}
% uses-lean: GlobalTheory.CuspRep.CuspSub, GlobalTheory.CuspRep.CG.IsRegAlg, GlobalTheory.ContraPair.PolarizedPair, GlobalTheory.WeilDeligne.wdOf, GlobalTheory.Pst.wdAt
\uses{fam:away-b}
Let $\mathrm{Rc}$ be a rec notion. \texttt{Row0b}$(\mathrm{Rc})$ is the statement: for every CM field $E$, algebraic closure $\Omega$, prime $l$, isomorphism $\iota\colon\overline{\Q}_l\cong\CC$, $a\ge1$, every cuspidal subrepresentation $\pi$ of the space of smooth functions on $\GL_a(\A_E)$ and every Hecke character $\chi$ of $E^+$ such that $\pi$ is regular algebraic and $(\pi,\chi)$ is a polarized pair ($\chi_v(-1)=(-1)^a$ at every real place $v$ of $E^+$, and $\pi$ twisted by $(\chi\circ N_{E/E^+})^{-1}$ is paired with $\pi$ through complex conjugation, the archive's \texttt{DualIsConjBy}), there are a homomorphism $r\colon G_E\to\GL_a(\overline{\Q}_l)$, an integer $\mathrm{wt}$ and multisets $\mathrm{HT}_{\tau}\subset\Z$ for the embeddings $\tau\colon E\to\overline{\Q}_l$ such that:
\begin{enumerate}
\item $r$ is continuous and semisimple;
\item at every finite $w$ with $l\notin w$, for every $\mathrm{rc}$ with $\mathrm{Rc}(\mathrm{rc})$ and every local component $\sigma$ of $\pi$ at $w$, \texttt{AwayB}$(\mathrm{WD}(r|_{G_w}),\mathrm{rc}(\sigma))$, where $\mathrm{WD}(r|_{G_w})$ is the archive's \texttt{wdOf};
\item $\mathrm{HT}_{\tau}$, viewed in $\CC$, is the multiset of $(a-1)/2-z$ for $z$ in the infinitesimal character of $\pi$ at $\iota\circ\tau$, and $\mathrm{HT}_{\tau\circ c}=\mathrm{wt}-\mathrm{HT}_{\tau}$;
\item for all local $l$-adic data at a place $w$ (a \texttt{LocPrime}, a \texttt{LocBridge} identifying the decomposition group with the Galois group of a finite extension of $\Q_l$, and compatible auxiliary data): if the local component $\sigma$ is spherical then $r|_{G_w}$ is crystalline and de Rham with Hodge--Tate jumps $\mathrm{HT}$; and for every $\mathrm{rc}$ with $\mathrm{Rc}(\mathrm{rc})$ such that $\mathrm{rc}(\sigma)$ has $N=0$ and either ($w$ is $\Q_g$-big for $a$ and $\mathrm{rc}(\sigma)$ is tame and square-free) or $\mathrm{rc}(\sigma)$ is abelian on commutators, $r|_{G_w}$ is de Rham and potentially semistable with jumps $\mathrm{HT}$, and \texttt{AwayB}$(\mathrm{WD}_{\mathrm{pst}}(r|_{G_w}),\mathrm{rc}(\sigma))$.
\end{enumerate}
\end{definition}

\begin{definition}[The stronger statement \texttt{Row0}]\label{fam:row0}
\lean{UpperTheory.GalStar.Away, UpperTheory.GalStar.Above, UpperTheory.GalStar.Props, UpperTheory.GalStar.Row0, UpperTheory.GalStar.IsQp}
\uses{fam:props-b, fam:away-b}
\texttt{Away}, \texttt{Above}, \texttt{Props} and \texttt{Row0} are \texttt{AwayB}, \texttt{AboveB}, \texttt{PropsB} and \texttt{Row0b} with ``$w$ is $\Q_g$-big for $n$'' replaced by ``$E_w\cong\Q_g$ for some prime $g$''.
\end{definition}

\begin{lemma}[\texttt{Row0} implies \texttt{Row0b}]\label{fam:row0-row0b}
\lean{UpperTheory.GalStar.Away.big, UpperTheory.GalStar.Props.big, UpperTheory.GalStar.Row0.row0b}
\uses{fam:row0, fam:props-b}
For every rec notion $\mathrm{Rc}$, \texttt{Row0}$(\mathrm{Rc})$ implies \texttt{Row0b}$(\mathrm{Rc})$.
\end{lemma}
\begin{proof}
A $\Q_g$-big place has $E_w\cong\Q_g$, so each clause of \texttt{Away} and \texttt{Above} gives the corresponding clause of \texttt{AwayB} and \texttt{AboveB}.
\end{proof}

\begin{definition}[The weak rec notion and \texttt{One.M.Row}]\label{fam:one-m-row}
\lean{UpperTheory.One.M.Row, UpperTheory.One.M.recT, UpperTheory.Gal.Int.recT}
\uses{fam:props-b, fam:laws-t}
The rec notion $\mathrm{recT}$ is \texttt{IsRecT} of Definition~\ref{fam:laws-t}. \texttt{One.M.Row} is \texttt{Row0b}$(\mathrm{recT})$.
\end{definition}

\begin{definition}[Existence of theta lifts]\label{fam:tlift-k}
\lean{UpperTheory.TP.TLiftK}
\texttt{TLiftK} is the statement: for every CM field $E$ with a good base, every $n\ge1$ and every hermitian space $X$ of dimension $n$ in the normalised form, there is $d_0$ such that for every $d\ge d_0$, every skew-hermitian space $W$ of dimension $d$ with scalar Gram matrix, all Hecke characters $\chi_X,\chi_W$ with prescribed restrictions and infinity types, and every automorphic representation $U$ of the unitary group of dimension $n$ with dominant weights $\lambda_w$ satisfying $d+m_W(w)\le 2\lambda_w(k)$ for all $w,k$, a theta lift $U'$ of $U$ to the unitary group of dimension $d$ exists.
\end{definition}

\begin{theorem}[Truncated local--global compatibility]\label{fam:gal-strong}
\lean{UpperTheory.Gal.Int.galoisLocalGlobalStrong_recParameter, UpperTheory.Gal.Int.galoisRepresentations_of_thetaLift, UpperTheory.Gal.Int.galoisLocalGlobalStrong_of_geometry_lift}
% uses-lean: UpperTheory.Gal.Int.Geo.GeoOne, UpperTheory.Gal.Int.Geo.geoOne_holds, UpperTheory.TP.K.tLiftK, UpperTheory.Gal.Int.galoisLocalGlobal_of_geometry_lift
\uses{fam:one-m-row, fam:tlift-k, fam:row0-row0b}
\texttt{Row0b}$(\mathrm{recT})$ holds. More precisely, assume \texttt{GeoOne} (a geometric realisation statement in weight one) and \texttt{TLiftK}; then \texttt{Row0b}$(\mathrm{recT})$; and both hypotheses are theorems of the archive.
\end{theorem}
\begin{proof}
\uses{fam:row0-row0b}
% uses-lean: UpperTheory.Gal.Int.coreAll_of, UpperTheory.Gal.Int.rest, UpperTheory.Gal.Int.Row.galoisLocalGlobal_of_localGlobalCompatibility
From \texttt{GeoOne} a seed statement in odd rank follows, and from the lacunary weight statement a congruence statement; with \texttt{TLiftK}, closed $\Phi$-functoriality, descent and limits of pseudo-characters this gives the core statement \texttt{CoreAll}. Local--global compatibility at spherical places, on unlinked supports and at Iwahori-fixed vectors then gives \texttt{Row0}$(\mathrm{recT})$, and Lemma~\ref{fam:row0-row0b} gives \texttt{Row0b}. \texttt{GeoOne} is \texttt{geoOne\_holds} and \texttt{TLiftK} is \texttt{tLiftK}. The proofs below this level lie in \arch{Development/AutomorphicArithmetic/Foundations/GaloisRepresentations/} and were not expanded.
\end{proof}

\begin{remark}[What \texttt{Row0b} omits, and what the quantifier over maps means]\label{fam:rem-row-strength}
\texttt{Row0b} asserts no comparison of monodromy away from $l$ beyond the case $N=0$, no purity, and at $w\nmid l$ an equality of Weil--Deligne representations only when the parameter is abelian on commutators; the tame case gives only characteristic polynomials on inertia, over $\Q_g$ with $g>16a^2$ or $a=1$. The archive proves the stronger \texttt{Row0}, which drops the bound on $g$, and then weakens it. The statement quantifies over every map with \texttt{IsRecT}. Since the weak laws fix such a map only on representations built from good cuspidals, the clauses have content only there: a value abelian on commutators of dimension at least $2$ cannot come from a cuspidal (cuspidals go to irreducibles), and the junk values of Lemma~\ref{fam:cusp-ex} are chosen, when $E_w\cong\Q_g$, not to be twists of tame representations. That no other map with \texttt{IsRecT} makes a clause fire on a non-good representation was not checked.
\end{remark}

\begin{remark}[Default Weil--Deligne parameters]\label{fam:rem-wd-default}
The parameter $\mathrm{WD}(r|_{G_w})$ at $w\nmid l$ (\texttt{wdOf}) is defined by choice and takes the trivial diagonal value when no parameter exists; that this branch is never taken would follow from Grothendieck's monodromy theorem, which was not located in the archive. The parameter at $w\mid l$ (\texttt{wdAt}) takes a trivial default when $r|_{G_w}$ is not potentially semistable; in clause 4 of Definition~\ref{fam:props-b} potential semistability is asserted together with \texttt{AwayB}, so that default is not reached there.
\end{remark}

\section{The literature slots at the top}



\begin{remark}[Two slots, one type]\label{fam:rem-row-dup}
The two fields named after Caraiani's local--global compatibility and after Galois representations of polarized pairs have the same type, the truncated statement of Definition~\ref{fam:one-m-row}, and are filled by the same theorem.
\end{remark}







\begin{remark}[Nothing is assumed at the top]\label{fam:rem-open-lit}
\texttt{OpenLiteratureStatements} is instantiated once, in Theorem~\ref{fin:every-character-of-row}, with closed theorems of the archive. The ``literature'' inputs are therefore claims proved inside the archive (not compiled here), of the reduced strength recorded in Remarks~\ref{fam:rem-not-llc} and~\ref{fam:rem-row-strength}.
\end{remark}
